\chapter{Introduction}


% \begin{equation}
%     \begin{tikzpicture}[baseline={(current bounding box.center)}]
%                 \begin{feynman}[small]
%                 \node [draw=red!80!black, circle, minimum size=0.5cm,
%            fill=red!15] (a) {\(F_j\)};
%                 \vertex [below right=1.7cm of a] (b) {\(y\)};
%                 \vertex [above right=1.7cm of a] (c) {\(x\)} ;
%                 \vertex [below left=1.7cm of a] (d) {\(2\)} ;
%                 \vertex [above left=1.7cm of a] (e) {\(1\)};
%                 \diagram*{
%                     (b) -- [] (a),
%                     (a) -- [] (c),
%                     (a) -- [] (d),
%                     (a) -- [] (e),
%                 };
%                 \end{feynman}
%                 \end{tikzpicture}
%                 \,
% 	\begin{tikzpicture}[baseline={(current bounding box.center)}]
% 	\draw[dashed,red,very thick] (0.0, -1.4) -- (0.0, 1.4);
% 	\end{tikzpicture}
% 	\,
%     \begin{tikzpicture}[baseline={(current bounding box.center)}]
%                 \begin{feynman}[small]
%                 \node [draw=blue!80!black, circle, minimum size=0.5cm,
%            fill=blue!15] (a) {\(\mathcal M\)};
%                 \vertex [below right=1.7cm of a] (b) {\(4\)};
%                 \vertex [above right=1.7cm of a] (c) {\(3\)} ;
%                 \vertex [below left=1.7cm of a] (d) {} ;
%                 \vertex [above left=1.7cm of a] (e) {};
%                 \diagram*{
%                     (a) -- [] (b),
%                     (a) -- [] (c),
%                     (a) -- [] (d),
%                     (a) -- [] (e),
%                 };
%                 \end{feynman}
%                 \end{tikzpicture}
% \end{equation}


1. frame inerente la tesi (il caso di fisica interessante che ha motivato lo studio)

2. stato dell'arte

3. i miei contributi in cui mi sono cimentato 
%(dai anche una panoramica più ampia) (risalta/motiva le tecniche usate, motiva l'approccio massless (nei problemi analizzati la parte UV è molto più importante per running e unitarietà, e qui nascono nuove simmetrie che rimarrebbero nascoste/oscurate nell'IR e poi menziona Nima per massive spinor helicity dicendo che alcuni pro rimangono nonostante il caso sia più complicato)) aggiungi anche i limiting factor, spiega perché le tecniche sono particolarmente efficaci nel contesto di EFT, semplificazioni clamorose (zeri, etc.), molta confidenza sui risultati, accesso agli insights della teoria.

4. spiega le tecniche e tools che ho imparato e messo a frutto

5. outline

\section{The Effective-Field-Theory Paradigm}

\subsection{Matching and Running}

\section{Synopsis of the Thesis}

The thesis is organized as follows.
Usa itemize.